%% 虚构示例：人物、学校、单位、项目、奖项、数字均不存在。请勿投递。
\documentclass[a4paper,fontset=none,zihao=5]{ctexart}
\usepackage{cnresume}
% \cnphoto{photo.jpg}
\hypersetup{pdftitle={程示例（虚构）},pdfauthor={程示例（虚构）},pdfsubject={虚构示例，请勿投递}}
\begin{document}
\cnheader{程示例（虚构）}{后端开发实习 · 杭州 · 2026.11可到岗}{%
  电话未提供\quad\href{mailto:campus@example.com}{campus@example.com}%
}
\cnnotice{虚构示例。人物、学校、单位、项目、竞赛、奖学金和数字均为编造，仅展示版式，请勿投递。}
\cnsection{教育背景}
\cnedu{2023.09-2027.06（预计）}{计算机科学与技术，本科在读}{示例虚构大学}{杭州}
\cnprofile{2027届（虚构）。虚构 GPA：3.6/4.0，排名：12/100。相关课程：数据库、操作系统、数据结构。}

\cnsection{实习经历}
\cnrole{2026.07-2026.08}{后端开发实习生}{示例虚构实验室}{杭州}
\begin{itemize}
  \item 在导师指导下维护教学用检索接口；个人完成输入校验与错误处理，未用于生产。
  \item 为接口补充虚构的20个测试用例；记录重复请求的问题和修复步骤。
\end{itemize}

\cnsection{项目经历}
\cnproject{示例课程预约系统}{2026.03-2026.06}{课程团队项目（虚构），负责接口与数据库设计}
\begin{itemize}
  \item 用 Python 和 SQLite 实现预约冲突检查，个人负责数据模型和查询接口。
  \item 在本机模拟虚构的100条预约，记录冲突场景；没有真实用户或线上性能数据。
\end{itemize}
\cnproject{示例日志分析练习}{2025.10-2025.12}{个人练习（虚构），仅使用合成数据}
\begin{itemize}
  \item 编写日志过滤脚本与使用说明，比较两种异常归类规则的局限。
\end{itemize}

\cnsection{在校经历 / 竞赛 / 奖学金}
\cnproject{示例校园编程赛}{2025.11}{虚构校级团队二等奖，负责测试与演示}
\begin{itemize}
  \item 组织队内任务拆分和一次复盘；团队成果不写成个人独立完成。
\end{itemize}
\cnprofile{虚构学习奖学金，示例虚构大学，2025.12。所有奖项只作排版示例。}

\cnsection{专业技能}
\cnskill{技术：}{Python、SQL、Git、单元测试；均为虚构展示}
\cnskill{语言：}{中文母语、英语课程阅读（虚构）}
\cnskill{实习安排：}{每周4天，持续3个月，课程时间另行协调（虚构）}
\end{document}
